\chapter{第5章：维数公式 习题}
\difficultykey

\section{基础题 \easy}

\begin{problem}{ch05-01}
\textbf{指数 $\mu$ 的计算。}
用公式
\[
\mu=[\SL_2(\Z):\GammaO(N)]=N\prod_{p\mid N}\left(1+\frac1p\right)
\]
计算 $N=1,2,3,4,6,12$ 时的 $\mu$。
\end{problem}
\begin{solution}
逐项代入即可：
\[
\begin{array}{c|c}
N & \mu \\ \hline
1 & 1 \\
2 & 2\left(1+\frac12\right)=3 \\
3 & 3\left(1+\frac13\right)=4 \\
4 & 4\left(1+\frac12\right)=6 \\
6 & 6\left(1+\frac12\right)\left(1+\frac13\right)=12 \\
12 & 12\left(1+\frac12\right)\left(1+\frac13\right)=24
\end{array}
\]
所以答案依次是 $1,3,4,6,12,24$。
\end{solution}

\begin{problem}{ch05-02}
\textbf{$N=1$ 时 $\Mk(\SL_2(\Z))$ 的维数。}
验证
\[
\dim \Mk(\SL_2(\Z))
\]
在 $k=4,6,8,10,12,14$ 时分别等于多少。
\end{problem}
\begin{solution}
对偶数 $k\ge 0$，标准公式是
\[
\dim \Mk(\SL_2(\Z))=
\begin{cases}
\lfloor k/12\rfloor,&k\equiv 2\pmod{12},\\
\lfloor k/12\rfloor+1,&k\not\equiv 2\pmod{12}.
\end{cases}
\]
于是
\[
\begin{array}{c|cccccc}
k & 4&6&8&10&12&14\\ \hline
\dim \Mk & 1&1&1&1&2&1
\end{array}
\]
这也可由 $\Mk(\SL_2(\Z))=\C[E_4,E_6]_k$ 直接看出：
$k=12$ 时基可取 $E_4^3,E_6^2$，其余所列权只有一个单项式基向量。
\end{solution}

\begin{problem}{ch05-03}
\textbf{唯一的权 $12$ 尖点形式。}
验证
\[
\dim S_{12}(\SL_2(\Z))=1,
\]
并说明 $\Delta$ 是唯一的归一化基向量。
\end{problem}
\begin{solution}
由上一题与关系
\[
S_{12}(\SL_2(\Z))=\Delta\cdot M_0(\SL_2(\Z))
\]
立刻得到
\[
\dim S_{12}(\SL_2(\Z))=\dim M_0(\SL_2(\Z))=1.
\]
更具体地说，任何权 $12$ 尖点形式都必须在唯一尖点 $\infty$ 处至少消失一次，因此必被判别式
\[
\Delta(q)=q\prod_{n\ge1}(1-q^n)^{24}
\]
整除。商是权 $0$ 模形式，只能是常数。所以 $S_{12}(\SL_2(\Z))=\C\Delta$。
归一化条件 $a_1=1$ 又唯一确定了基向量，因此归一化基向量就是 $\Delta$。
\end{solution}

\begin{problem}{ch05-04}
\textbf{$\dim S_2(\GammaO(11))$ 与亏格。}
用公式
\[
\dim S_2(\GammaO(N))=g_0(N)=1+\frac{\mu}{12}-\frac{\nu_2}{4}-\frac{\nu_3}{3}-\frac{\nu_\infty}{2}
\]
计算 $\dim S_2(\GammaO(11))$，并与 $g_0(11)$ 比较。
\end{problem}
\begin{solution}
对 $N=11$：
\[
\mu=11\left(1+\frac1{11}\right)=12.
\]
又因为 $11\equiv 3\pmod4$，故 $\left(\frac{-1}{11}\right)=-1$，从而 $\nu_2=1+(-1)=0$。
同理 $11\equiv 2\pmod3$，故 $\left(\frac{-3}{11}\right)=-1$，从而 $\nu_3=0$。
尖点数对素数 $p$ 为 $2$，所以 $\nu_\infty=2$。
于是
\[
\dim S_2(\GammaO(11))=1+\frac{12}{12}-0-0-1=1.
\]
因此 $g_0(11)=1$，与 $\dim S_2(\GammaO(11))$ 完全一致。
\end{solution}

\begin{problem}{ch05-05}
\textbf{$\dim S_2(\GammaO(37))$。}
计算 $\dim S_2(\GammaO(37))$，并解释为什么这意味着 $J_0(37)$ 分解出两个彼此独立的一维椭圆曲线商。
\end{problem}
\begin{solution}
对 $N=37$：
\[
\mu=37\left(1+\frac1{37}\right)=38,
\qquad \nu_\infty=2.
\]
因为 $37\equiv 1\pmod4$，所以 $\nu_2=1+1=2$；
又 $37\equiv 1\pmod3$，所以 $\nu_3=1+1=2$。
故
\[
\dim S_2(\GammaO(37))=1+\frac{38}{12}-\frac24-\frac23-1=2.
\]
因此 $J_0(37)$ 是二维 Abel 簇。Hecke 代数在该二维空间上可同时对角化，
而实际上这两个本征子空间都由有理系数新形式生成，所以对应两个一维 Abel 子簇，也就是两条椭圆曲线商。
换言之，$J_0(37)$ 与两条导子 $37$ 的椭圆曲线的乘积同源。
\end{solution}

\begin{problem}{ch05-06}
\textbf{计算 $\nu_2$。}
对素数 $p$，若 $p\neq 2$，则
\[
\nu_2=1+\left(\frac{-1}{p}\right).
\]
计算 $p=5,7,11,13$ 时的 $\nu_2$。
\end{problem}
\begin{solution}
利用
\[
\left(\frac{-1}{p}\right)=(-1)^{(p-1)/2}
\]
即可：
\[
\begin{array}{c|c|c}
p & \left(\frac{-1}{p}\right) & \nu_2 \\ \hline
5 & 1 & 2 \\
7 & -1 & 0 \\
11 & -1 & 0 \\
13 & 1 & 2
\end{array}
\]
因此答案是
\[
\nu_2(5)=2,\quad \nu_2(7)=0,\quad \nu_2(11)=0,\quad \nu_2(13)=2.
\]
\end{solution}

\begin{problem}{ch05-07}
\textbf{计算 $\nu_3$。}
对素数 $p\neq 3$，有
\[
\nu_3=1+\left(\frac{-3}{p}\right).
\]
计算 $p=7,13,19$ 时的 $\nu_3$。
\end{problem}
\begin{solution}
逐个判断 $-3$ 是否为模 $p$ 的平方剩余：
\begin{itemize}
  \item $p=7$ 时，$-3\equiv 4\pmod7$，而 $4$ 是平方，故 $\left(\frac{-3}{7}\right)=1$；
  \item $p=13$ 时，$-3\equiv 10\pmod{13}$，也可由二次互反律算出 $\left(\frac{-3}{13}\right)=1$；
  \item $p=19$ 时，$-3\equiv 16\pmod{19}$，显然也是平方，故符号为 $1$。
\end{itemize}
所以三者都满足
\[
\nu_3=1+1=2.
\]
即
\[
\nu_3(7)=\nu_3(13)=\nu_3(19)=2.
\]
\end{solution}

\begin{problem}{ch05-08}
\textbf{$N=1,\dots,12$ 时的 $\dim S_2(\GammaO(N))$。}
利用亏格公式列出 $N=1,\dots,12$ 时 $\dim S_2(\GammaO(N))$ 的表格。
\end{problem}
\begin{solution}
计算结果如下：
\[
\begin{array}{c|cccccccccccc}
N&1&2&3&4&5&6&7&8&9&10&11&12\\ \hline
\dim S_2(\GammaO(N))&0&0&0&0&0&0&0&0&0&0&1&0
\end{array}
\]
其中唯一非零的是 $N=11$。这也与经典事实一致：
$X_0(N)$ 在 $N\le 10$ 以及 $N=12$ 时皆为亏格 $0$，而 $X_0(11)$ 是最小的亏格 $1$ 模曲线。
\end{solution}

\begin{problem}{ch05-09}
\textbf{Eisenstein 部分与尖点数。}
验证对偶权 $k\ge 2$，有
\[
\dim \Mk(\GammaO(N))=\dim \Sk(\GammaO(N))+\nu_\infty.
\]
并解释这条公式的含义。
\end{problem}
\begin{solution}
对偶权 $k\ge 2$，$\Mk(\GammaO(N))$ 分解为尖点部分和 Eisenstein 部分：
\[
\Mk(\GammaO(N))=\Sk(\GammaO(N))\oplus \mathcal E_k(\GammaO(N)).
\]
对 $\GammaO(N)$ 的每个尖点都可构造一个“以该尖点为常数项坐标”的 Eisenstein 级数；
这些级数张成的空间维数正好等于尖点数 $\nu_\infty$。
因此
\[
\dim \mathcal E_k(\GammaO(N))=\nu_\infty,
\]
从而得到
\[
\dim \Mk(\GammaO(N))=\dim \Sk(\GammaO(N))+\nu_\infty.
\]
含义是：非尖点模形式无非是在尖点形式之外，再加上“给每个尖点一个常数项自由度”。
\end{solution}

\begin{problem}{ch05-10}
\textbf{$N=11$ 时的 Eisenstein 维数。}
用上一题的结论验证对任意偶权 $k\ge 2$，有
\[
\dim M_k(\GammaO(11))-\dim S_k(\GammaO(11))=2.
\]
\end{problem}
\begin{solution}
因为 $11$ 是素数，$X_0(11)$ 只有两个尖点 $\infty$ 与 $0$，故
\[
\nu_\infty(\GammaO(11))=2.
\]
由上一题的一般结论立刻得到
\[
\dim M_k(\GammaO(11))-\dim S_k(\GammaO(11))=\nu_\infty=2.
\]
这说明每个偶权 $k\ge 2$ 都恰有两个线性独立的 Eisenstein 级数，一个对应尖点 $\infty$，一个对应尖点 $0$。
\end{solution}

\section{进阶题 \medium}

\begin{problem}{ch05-11}
\textbf{$\dim M_k(\SL_2(\Z))$ 的完整公式。}
从模曲线 $X(1)$ 上的 Riemann--Roch 出发，推导
\[
\dim M_k(\SL_2(\Z))=
\begin{cases}
0,&k<0\text{ 或 }k\text{ 奇},\\
1,&k=0,\\
\lfloor k/12\rfloor,&k\equiv 2\pmod{12},\ k\ge2,\\
\lfloor k/12\rfloor+1,&k\not\equiv 2\pmod{12},\ k\ge2.
\end{cases}
\]
\end{problem}
\begin{solution}
模曲线 $X(1)$ 亏格为 $0$，但在椭圆点 $i,\rho$ 与尖点 $\infty$ 处有 orbifold 修正。
权 $k$ 模形式对应于一个线丛 $\mathcal L^{\otimes k}$ 的全局截面，其“有效次数”是
\[
\deg(\mathcal L^{\otimes k})=\frac{k}{12}.
\]
更准确地说，一个权 $k$ 模形式在 $i$ 处最多允许极点阶 $\lfloor k/4\rfloor-k/4$ 的补偿，在 $\rho$ 处最多允许
$\lfloor k/3\rfloor-k/3$ 的补偿，在尖点处则由 $q$-展开的非负性控制。

因此可把维数理解为满足
\[
4a+6b=k
\]
的非负整数解 $(a,b)$ 的个数，因为
\[
M_*(\SL_2(\Z))=\C[E_4,E_6].
\]
这立刻给出：
\begin{itemize}
  \item 若 $k$ 为奇数，则无解，所以维数为 $0$；
  \item 若 $k\ge 0$ 为偶数，则解的个数等于上式中的分段函数。
\end{itemize}
具体看模 $12$ 的余数：
\begin{itemize}
  \item 当 $k\equiv 0,4,6,8,10\pmod{12}$ 时，总能先取一个显然的基解，再每增加 $12$ 贡献一个解，所以维数是 $\lfloor k/12\rfloor+1$；
  \item 当 $k\equiv 2\pmod{12}$ 时，最小的非负解不存在，只能从 $k-14,k-26,\dots$ 往上平移，故少一个，维数是 $\lfloor k/12\rfloor$。
\end{itemize}
这与 Riemann--Roch 观点完全一致：在 $\mathbb P^1$ 上，线丛次数加一给出维数，而椭圆点的分数修正正好造成 $k\equiv2\pmod{12}$ 时的“少一维”现象。
\end{solution}

\begin{problem}{ch05-12}
\textbf{尖点数公式。}
证明
\[
\nu_\infty(\GammaO(N))=\sum_{d\mid N}\phi\bigl(\gcd(d,N/d)\bigr).
\]
\end{problem}
\begin{solution}
$\GammaO(N)$ 在 $\mathbb P^1(\Q)$ 上的轨道由既约分数 $a/c$ 的分母 $c$ 的约数类决定。
更具体地说，每个尖点都可表示为某个 $1/d$，其中 $d\mid N$。

固定 $d\mid N$。与 $1/d$ 同分母层的尖点，其稳定子中的平移宽度只取决于
\[
\gcd(d,N/d).
\]
当分母层固定时，互不同余的分子类由
\[
(\Z/\gcd(d,N/d)\Z)^\times
\]
参数化，因此该层恰有
\[
\phi(\gcd(d,N/d))
\]
个尖点。

对所有 $d\mid N$ 求和，就得到
\[
\nu_\infty(\GammaO(N))=\sum_{d\mid N}\phi\bigl(\gcd(d,N/d)\bigr).
\]
例如 $N=p$ 为素数时，约数只有 $1,p$，两项都给出 $\phi(1)=1$，故 $\nu_\infty=2$。
\end{solution}

\begin{problem}{ch05-13}
\textbf{Nebentypus 情形的维数公式。}
对带特征的空间 $M_k(\GammaO(N),\chi)$，说明维数公式如何修改：
哪些项与 $\chi$ 有关，哪些项保持不变？
\end{problem}
\begin{solution}
带特征时，模形式对应的不再是平凡线丛上的截面，而是带有扭曲自守因子的线丛截面。
因此维数公式的整体形状仍由 Riemann--Roch 给出，但局部修正要乘上 $\chi$ 在稳定子上的取值条件。

保持不变的部分有：
\begin{itemize}
  \item 指数 $\mu=[\SL_2(\Z):\GammaO(N)]$；
  \item 模曲线的几何亏格；
  \item 尖点的组合结构与宽度。
\end{itemize}
会随 $\chi$ 变化的部分有：
\begin{itemize}
  \item 椭圆点修正：若 $\chi$ 在阶 $2$ 或阶 $3$ 稳定子上的取值与权 $k$ 的自守因子不相容，则相应椭圆点根本不贡献截面；
  \item Eisenstein 部分：只有当某个尖点处的常数项与 $\chi$ 兼容时，该尖点才贡献一维 Eisenstein 方向，所以“尖点数”要替换成“与 $\chi$ 相容的尖点数”；
  \item 权 $1$ 与某些低权时还会出现额外的特殊现象。
\end{itemize}
因此可以把无特征公式理解成：把每个局部稳定子处的贡献，按 $\chi$-本征条件筛选一遍。形式上仍是
\[
\dim M_k(\GammaO(N),\chi)=\text{主项 }\frac{(k-1)\mu}{12}+\text{局部修正},
\]
只是局部修正中的椭圆点与尖点项都要按 $\chi$ 的值进行删减。
\end{solution}

\begin{problem}{ch05-14}
\textbf{素数水平的新形式维数。}
对 $N=p$ 为素数，解释为什么
\[
S_k(\GammaO(p))=S_k^{\mathrm{new}}(\GammaO(p))\oplus S_k^{\mathrm{old}}(\GammaO(p)),
\]
且在权 $2$ 时 $S_2^{\mathrm{old}}(\GammaO(p))=0$。据此计算
\[
\dim S_2^{\mathrm{new}}(\GammaO(p))
\]
对 $p=11,17,19,23$ 的值。
\end{problem}
\begin{solution}
一般地，素数水平的旧形式都来自 level $1$，通过
\[
f(\tau)\mapsto f(\tau),\qquad f(\tau)\mapsto f(p\tau)
\]
两个退化嵌入得到，因此旧空间同构于 $S_k(\SL_2(\Z))^{\oplus 2}$。
但在权 $2$ 时
\[
S_2(\SL_2(\Z))=0,
\]
所以
\[
S_2^{\mathrm{old}}(\GammaO(p))=0,
\qquad
S_2^{\mathrm{new}}(\GammaO(p))=S_2(\GammaO(p)).
\]
由亏格公式可得
\[
\dim S_2^{\mathrm{new}}(\GammaO(11))=1,
\quad
\dim S_2^{\mathrm{new}}(\GammaO(17))=1,
\quad
\dim S_2^{\mathrm{new}}(\GammaO(19))=1,
\quad
\dim S_2^{\mathrm{new}}(\GammaO(23))=2.
\]
因此前三个水平各对应一条新椭圆曲线，$N=23$ 则对应两个独立的新方向。
\end{solution}

\begin{problem}{ch05-15}
\textbf{较大水平的计算。}
计算
\[
\dim S_2(\GammaO(N))
\]
在 $N=50,100,200$ 时的值，并比较它们与线性主项的关系。
\end{problem}
\begin{solution}
先算组合量：
\[
\begin{array}{c|ccc}
N & 50 & 100 & 200 \\ \hline
\mu & 90 & 180 & 360 \\
\nu_2 & 2 & 0 & 0 \\
\nu_3 & 0 & 0 & 0 \\
\nu_\infty & 12 & 18 & 24
\end{array}
\]
故
\[
\dim S_2(\GammaO(50))=1+\frac{90}{12}-\frac24-6=2,
\]
\[
\dim S_2(\GammaO(100))=1+\frac{180}{12}-0-9=7,
\]
\[
\dim S_2(\GammaO(200))=1+\frac{360}{12}-0-12=19.
\]
所以答案是
\[
2,\ 7,\ 19.
\]
这说明维数随 $N$ 线性增长。更精确地说，这三个 $N$ 都只有素因子 $2,5$，于是
\[
\mu=N\left(1+\frac12\right)\left(1+\frac15\right)=\frac{9N}{5},
\]
故主项其实是 $\mu/12=3N/20$，而不是严格的 $N/12$。因此应把“$N/12$”理解成粗略量级；
真正的主项由 $\mu/12$ 控制。
\end{solution}

\begin{problem}{ch05-16}
\textbf{何时 $\dim S_2(\GammaO(N))=0$？}
证明
\[
\dim S_2(\GammaO(N))=0
\]
当且仅当 $X_0(N)$ 亏格为 $0$，并写出全部这样的 $N$。
\end{problem}
\begin{solution}
因为
\[
\dim S_2(\GammaO(N))=g_0(N),
\]
所以 $\dim S_2(\GammaO(N))=0$ 与 $g_0(N)=0$ 完全等价。
而经典结果指出
\[
X_0(N)\text{ 亏格为 }0
\iff
N\in\{1,2,3,4,5,6,7,8,9,10,12,13,16,18,25\}.
\]
因此全部满足条件的 $N$ 恰是上述有限集合。
这也解释了为什么从 $N=11$ 开始就首次出现非平凡的权 $2$ 尖点形式。
\end{solution}

\begin{problem}{ch05-17}
\textbf{$g=1$ 时的 Riemann--Roch。}
在 $X_0(11)$ 上验证
\[
\ell(D)-\ell(K-D)=\deg D-g+1
\]
在 $g=1$ 的情形成立。
\end{problem}
\begin{solution}
对 $X_0(11)$，我们知道 $g=1$，而椭圆曲线上的典范因子满足
\[
K\sim 0.
\]
因此 Riemann--Roch 变成
\[
\ell(D)-\ell(-D)=\deg D.
\]
若 $\deg D>0$，则 $-D$ 的次数为负，不可能有非零全纯函数以 $-D$ 为极点上界，所以
\[
\ell(-D)=0,
\qquad \ell(D)=\deg D.
\]
例如取 $D=(P)$，则 $\deg D=1$，得到 $\ell(P)=1$；这正说明椭圆曲线上次数 $1$ 的除子类给出一维函数空间。
若 $D=0$，则
\[
\ell(0)-\ell(0)=0,
\]
自然成立，其中 $\ell(0)=1$ 是常数函数空间。
所以在 $g=1$ 时，Riemann--Roch 的内容就是：正次数除子的截面维数恰好等于其次数。
\end{solution}

\begin{problem}{ch05-18}
\textbf{权 $1$ 的特殊性。}
陈述 Selberg--Stark 思想：$\dim S_1(\GammaO(N))$ 与某些 Artin $L$-函数在 $s=1$ 的正则性有关。并给出
\[
\dim S_1(\GammaO(12)),\qquad \dim S_1(\GammaO(23))
\]
的值。
\end{problem}
\begin{solution}
权 $1$ 不能简单套用高权维数公式，因为 Eichler--Shimura 对权 $1$ 不再直接给出几何实现。
Selberg--Stark 的观点是：权 $1$ 尖点形式对应二维复 Artin 表示，
于是其维数与相应 Artin $L$-函数在 $s=1$ 的解析性质紧密相关。

具体数值上，查经典表或 LMFDB 可得
\[
\dim S_1(\GammaO(12))=0,
\qquad
\dim S_1(\GammaO(23))=2.
\]
其中 $N=23$ 的情形来自一个二面体型权 $1$ 新形式及其 Galois 共轭，因此总维数为 $2$。
这正说明权 $1$ 空间虽小，却承载了深刻的 Galois 信息。
\end{solution}

\begin{problem}{ch05-19}
\textbf{乘以 $\Delta$ 的同构。}
证明乘以判别式给出同构
\[
\Mk(\SL_2(\Z))\xrightarrow{\ \cdot\Delta\ } S_{k+12}(\SL_2(\Z)).
\]
据此推出
\[
\dim S_{k+12}(\SL_2(\Z))=\dim M_k(\SL_2(\Z)).
\]
\end{problem}
\begin{solution}
映射
\[
f\longmapsto f\Delta
\]
显然把权 $k$ 模形式送到权 $k+12$ 尖点形式，因为 $\Delta$ 本身是权 $12$ 尖点形式，且在 $\infty$ 处恰有一次零点。

它是单射，因为 $\Delta\neq 0$。下面证满射。
若 $g\in S_{k+12}(\SL_2(\Z))$，则 $g$ 在唯一尖点 $\infty$ 处至少消失一次；
而 $\Delta$ 也恰好在该尖点消失一次，且在上半平面无零点，所以
\[
\frac{g}{\Delta}
\]
在上半平面全纯，在尖点处仍全纯，因此是权 $k$ 模形式。
于是 $g=f\Delta$，其中 $f\in M_k(\SL_2(\Z))$。

故乘以 $\Delta$ 是双射，从而
\[
S_{k+12}(\SL_2(\Z))\cong M_k(\SL_2(\Z)),
\qquad
\dim S_{k+12}=\dim M_k.
\]
\end{solution}

\begin{problem}{ch05-20}
\textbf{$\SL_2(\Z)$ 的 Eisenstein 维数。}
推导对偶权 $k\ge2$，有
\[
\dim \Mk(\SL_2(\Z))=\dim \Sk(\SL_2(\Z))+1.
\]
\end{problem}
\begin{solution}
因为 $X(1)$ 只有一个尖点 $\infty$，所以 Eisenstein 子空间只有一维。
具体地，任一偶权 $k\ge 4$ 都有经典 Eisenstein 级数 $E_k$，它张成全部非尖点方向。
因此
\[
M_k(\SL_2(\Z))=S_k(\SL_2(\Z))\oplus \C E_k.
\]
于是
\[
\dim M_k(\SL_2(\Z))=\dim S_k(\SL_2(\Z))+1.
\]
若要从上一题推出，也可写成
\[
\dim S_k=\dim M_{k-12},
\]
再与 $M_k$ 的显式维数公式比较，差恰好也是 $1$。
\end{solution}

\begin{problem}{ch05-21}
\textbf{一般 $\GammaO(N)$ 的直和分解。}
证明对偶权 $k\ge2$，有
\[
\Mk(\GammaO(N))=\Sk(\GammaO(N))\oplus \mathcal E_k(\GammaO(N)),
\qquad
\dim \mathcal E_k(\GammaO(N))=\nu_\infty.
\]
\end{problem}
\begin{solution}
先看交为零。若 $f$ 同时属于 $S_k$ 与 $\mathcal E_k$，那么它在每个尖点的常数项都为零；
但 Eisenstein 级数空间正是由“允许非零常数项”的模形式组成，常数项全为零就只能落回尖点空间的零元，故交为零。

再看张成性。任取 $f\in M_k(\GammaO(N))$，记其在各尖点的常数项为一组数
\[
(c_1,\dots,c_{\nu_\infty}).
\]
对每个尖点可选择一个 Eisenstein 级数 $E_j$，使它在第 $j$ 个尖点的常数项为 $1$，在其余尖点按某固定归一化给出。
这些常数项矩阵可证明可逆，因此能找到线性组合
\[
E=\sum a_jE_j
\]
使得 $f-E$ 在所有尖点的常数项都为零，于是 $f-E\in S_k(\GammaO(N))$。
故
\[
M_k=S_k+\mathcal E_k.
\]
由交零可知这是直和。

最后，由“每个尖点给一个独立常数项自由度”可知
\[
\dim \mathcal E_k(\GammaO(N))=\nu_\infty.
\]
\end{solution}

\begin{problem}{ch05-22}
\textbf{一维尖点空间与椭圆曲线。}
证明若
\[
\dim S_2(\GammaO(N))=1,
\]
则 $J_0(N)$ 是一条椭圆曲线。
\end{problem}
\begin{solution}
Jacobi 簇的维数满足
\[
\dim J_0(N)=g_0(N)=\dim S_2(\GammaO(N)).
\]
若右边等于 $1$，则 $J_0(N)$ 是一维 Abel 簇。
而一维 Abel 簇的定义正是一条带原点的椭圆曲线。
因此 $J_0(N)$ 本身就是一条椭圆曲线。

换句话说，权 $2$ 尖点空间只有一个 Hecke 本征方向时，整个 Jacobian 已无法再分解成更高维的 Abel 成分。
\end{solution}

\begin{problem}{ch05-23}
\textbf{选做：权 $3/2$ 的维数公式。}
简述半整数权 $3/2$ 模形式的维数公式与 Shimura 对应之间的关系。
\end{problem}
\begin{solution}
半整数权情形要把群换成某个双覆盖并引入 theta 乘子系统，因此维数公式的主项仍来自 Riemann--Roch，
但局部修正不再只取决于椭圆点与尖点，还取决于乘子系统在稳定子上的值。

对权 $3/2$ 尖点形式，Kohnen plus 空间是最重要的子空间。Shimura 对应给出一个 Hecke-equivariant 的映射
\[
S_{3/2}^{+}(4N)\longrightarrow S_2(N),
\]
把半整数权形式与权 $2$ 新形式联系起来。
因此，权 $3/2$ 的维数往往可理解为：权 $2$ 空间的维数，再加上一些 theta 系列造成的补偿项。
这也是为什么半整数权维数公式看起来更复杂，但其核心仍受权 $2$ 理论控制。
\end{solution}

\begin{problem}{ch05-24}
\textbf{$N\le 100$ 时哪几个水平满足 $\dim S_2(\GammaO(N))=1$？}
用维数表格找出全部满足
\[
\dim S_2(\GammaO(N))=1
\]
且 $1\le N\le 100$ 的 $N$。
\end{problem}
\begin{solution}
因为 $\dim S_2(\GammaO(N))=g_0(N)$，所以只需找出亏格为 $1$ 的水平。
计算可得
\[
N\in\{11,14,15,17,19,20,21,24,27,32,36,49\}.
\]
这些正是 $N\le100$ 时使得 $J_0(N)$ 成为椭圆曲线的全部水平。
在模性定理与 Cremona 表中，它们都具有特殊地位。
\end{solution}

\begin{problem}{ch05-25}
\textbf{权的最终线性与模 $12$ 周期。}
证明对固定 $N$，$\dim M_k(\GammaO(N))$ 对足够大的偶权 $k$，在每个模 $12$ 的剩余类上都是关于 $k$ 的仿射线性函数。
\end{problem}
\begin{solution}
把 $M_k(\GammaO(N))$ 视为模曲线上的某个高次线丛截面空间。由 Riemann--Roch，
对 $k\gg 0$ 有
\[
\dim M_k(\GammaO(N))=\deg(\mathcal L_k)-g_0(N)+1,
\]
其中 $\deg(\mathcal L_k)$ 对 $k$ 线性，主系数为 $\mu/12$。

另一方面，椭圆点修正只取决于
\[
\left\lfloor \frac{k}{4}\right\rfloor,
\qquad
\left\lfloor \frac{k}{3}\right\rfloor,
\]
所以它们仅依赖于 $k$ 模 $12$ 的余数类。
于是当 $k$ 足够大时，在每个固定的剩余类上都有
\[
\dim M_k(\GammaO(N))=\frac{\mu}{12}k+c_r,
\qquad k\equiv r\pmod{12},
\]
其中常数 $c_r$ 只与 $r$ 有关。

等价地说：维数函数整体上不是严格周期的，但“去掉线性主项后”只剩下一个模 $12$ 周期的有界振荡项。
\end{solution}

\section{挑战题 \hard}

\begin{problem}{ch05-26}
\textbf{Eichler--Selberg 迹公式的最小检验。}
陈述
\[
\Tr(T_n\mid S_k(\SL_2(\Z)))
\]
的 Eichler--Selberg 迹公式的结构，并对 $k=12$、$n=p$ 为素数时解释为什么
\[
\Tr(T_p)=\tau(p).
\]
\end{problem}
\begin{solution}
Eichler--Selberg 迹公式把 $\Tr(T_n)$ 写成三类贡献之和：
\begin{itemize}
  \item 抛物项；
  \item 椭圆项（涉及二次型类数）；
  \item 双曲项。
\end{itemize}
这是 Hecke 算子在模形式空间上的“显式求迹公式”。

对 $k=12$，我们知道
\[
\dim S_{12}(\SL_2(\Z))=1,
\]
所以 $S_{12}(\SL_2(\Z))$ 由唯一归一化特征形式 $\Delta$ 张成。Hecke 算子在一维空间上只能按标量作用，而该标量恰是 $\Delta$ 的第 $p$ 个 Hecke 本征值，即 Ramanujan 的 $\tau(p)$。
因此
\[
\Tr(T_p\mid S_{12}(\SL_2(\Z)))=\tau(p).
\]
换言之，在一维情形下，迹公式自动退化为“特征值本身”。
\end{solution}

\begin{problem}{ch05-27}
\textbf{Verdier 对偶与权的互补。}
解释为什么可以把
\[
\mathbb D(S_k(\Gamma))\cong M_{2-k}(\Gamma)^*
\]
看成“权 $k$ 与权 $2-k$ 互补”的几何表达。
\end{problem}
\begin{solution}
在曲线上，Verdier 对偶把一个截面空间转成与典范层有关的对偶截面空间。
模形式权 $k$ 对应于某个自守线丛 $\omega^{\otimes k}$，而对偶时会多出一个典范层因子 $\omega^{\otimes 2}$。
因此对偶权自然从 $k$ 变成
\[
2-k.
\]

解析上，这与 Petersson 配对及 Serre 对偶是一回事：
尖点形式因为在尖点处消失，恰好与低权弱模形式形成配对。
所以公式
\[
\mathbb D(S_k(\Gamma))\cong M_{2-k}(\Gamma)^*
\]
并不是神秘的新对象，而正是“模曲线上的 Serre 对偶”在自守线丛语言中的写法。
\end{solution}

\begin{problem}{ch05-28}
\textbf{模符号与有理结构。}
定义
\[
\{r,s\}=\int_r^s f(\tau)\,(2\pi i\,d\tau),\qquad r,s\in\mathbb P^1(\Q).
\]
说明模符号为什么给出 $S_2(\GammaO(N))$ 的有理结构。
\end{problem}
\begin{solution}
对 $f\in S_2(\GammaO(N))$，积分
\[
\int_r^s f(\tau)\,d\tau
\]
只依赖于从 $r$ 到 $s$ 的同调类，因此它定义了
\[
H_1\bigl(X_0(N),\{\text{cusps}\};\Z\bigr)\longrightarrow \C
\]
上的线性泛函。Manin 发现，这个相对同调群可由有限个模符号生成，并满足显式整数关系。

于是 $S_2(\GammaO(N))$ 的对偶可嵌入到一个有限生成的整数模张量 $\Q$ 后得到的向量空间中，
从而自然获得一个 $\Q$-结构。
换言之，Hecke 算子对模符号的整数作用让权 $2$ 尖点形式不只是一个复向量空间，
而是一个带 Hecke 兼容有理结构的对象。
\end{solution}

\begin{problem}{ch05-29}
\textbf{共轭子群下维数不变。}
证明若 $\Gamma' = g\Gamma g^{-1}$，其中 $g\in \SL_2(\Z)$，则
\[
\dim M_k(\Gamma)=\dim M_k(\Gamma').
\]
\end{problem}
\begin{solution}
定义算子
\[
(f\mid_k g)(\tau)=(c\tau+d)^{-k}f\!\left(\frac{a\tau+b}{c\tau+d}\right),
\qquad
g=\begin{pmatrix}a&b\\ c&d\end{pmatrix}.
\]
若 $f\in M_k(\Gamma)$，则对任意 $\gamma'\in \Gamma'=g\Gamma g^{-1}$，有
\[
(f\mid_k g)\mid_k \gamma'=f\mid_k(g\gamma')=f\mid_k(\gamma g)=f\mid_k g.
\]
所以 $f\mapsto f\mid_k g$ 给出
\[
M_k(\Gamma)\xrightarrow{\sim} M_k(\Gamma').
\]
这是线性同构，因此维数相等。

本质上，共轭只是在上半平面上换坐标；模曲线的几何对象并未改变。
而对非共轭子群，这种几何同构不再存在，所以维数一般可能不同。
\end{solution}

\begin{problem}{ch05-30}
\textbf{有理 Hecke 特征值的例子。}
给出两个水平 $N$ 的例子，使得 $S_k(\GammaO(N))$ 中全部 Hecke 特征值都落在 $\Q$ 中，并说明验证方法。
\end{problem}
\begin{solution}
最简单的例子是
\[
N=11,\qquad N=37,
\]
都取权 $k=2$。

\textbf{例 1：$N=11$。}
$S_2(\GammaO(11))$ 只有一维，由唯一归一化新形式
\[
f_{11}=q-2q^2-q^3+2q^4+q^5+\cdots
\]
张成。所有 Hecke 算子都在一维空间上按有理数作用，所以全部特征值都在 $\Q$ 中。

\textbf{例 2：$N=37$。}
$S_2(\GammaO(37))$ 二维，但它分解成两个一维有理新形式方向，分别对应两条导子 $37$ 的椭圆曲线。
因此每个 Hecke 本征值都就是某条椭圆曲线的 Frobenius 迹 $a_p(E)\in\Z$，仍落在 $\Q$ 中。

验证方法很统一：
\begin{itemize}
  \item 先把空间分解成新形式方向；
  \item 再检查每个新形式的 Hecke 域 $K_f$ 是否等于 $\Q$；
  \item 若 $K_f=\Q$，则其全部 Fourier 系数与 Hecke 特征值都为有理数。
\end{itemize}
\end{solution}
